\documentclass[conference]{IEEEtran}
\IEEEoverridecommandlockouts
\usepackage{cite}
\usepackage{amsmath,amssymb,amsfonts}
\usepackage{graphicx}
\usepackage{textcomp}
\usepackage{xcolor}
\usepackage{booktabs}
\usepackage{array}
\usepackage{tabularx}
\usepackage[hidelinks]{hyperref}
\graphicspath{{figures/}{../../../paper/generated/}}
\setlength{\textfloatsep}{7pt plus 2pt minus 2pt}
\setlength{\floatsep}{7pt plus 2pt minus 2pt}
\setlength{\intextsep}{7pt plus 2pt minus 2pt}
\renewcommand{\arraystretch}{0.98}
\newcommand{\AuditTotalRows}{1257}
\newcommand{\AuditValidationRows}{140}
\newcommand{\AuditValidationExcluded}{4}
\newcommand{\AuditTestCandidateRows}{343}
\newcommand{\AuditTestExcluded}{10}
\newcommand{\AuditTestRows}{333}
\newcommand{\AuditTailRows}{73}
\newcommand{\AuditSelectedModel}{ExtraTrees}
\newcommand{\AuditSelectedFamily}{weather only}
\newcommand{\AuditSelectedRtwo}{-0.014}
\newcommand{\AuditSelectedRMSE}{0.669}
\newcommand{\AuditSelectedMAE}{0.474}
\newcommand{\AuditSelectedDeltaRMSE}{-0.005}
\newcommand{\AuditSelectedDeltaHigh}{0.009}
\newcommand{\AuditGateACI}{[-0.019, 0.009]}
\newcommand{\AuditGateBOneDeltaRMSE}{-0.012}
\newcommand{\AuditGateBOneCI}{[-0.029, 0.002]}
\newcommand{\AuditGateBTwoCI}{[-0.019, 0.009]}
\newcommand{\AuditTailGate}{FAIL}
\newcommand{\AuditPositiveStates}{4/12}
\newcommand{\AuditLeakageJaccard}{0.605}
\newcommand{\AuditLeakageRank}{0.942}

\newcommand{\SyntheticPolicyFalsePermissions}{171/240}
\newcommand{\SyntheticPolicyFalsePermissionRate}{71.2\%}
\newcommand{\SyntheticPolicyFalseAbstentions}{20/180}
\newcommand{\SyntheticPolicyFalseAbstentionRate}{11.1\%}
\newcommand{\SyntheticPolicyPermissionRate}{78.8\%}
\newcommand{\SyntheticPolicySensitivity}{88.9\%}
\newcommand{\SyntheticPolicySpecificity}{28.7\%}


\title{Auditing Weather-Feature Reliance in Detrended Crop-Yield Models}
\author{\IEEEauthorblockN{Anonymous Submission}}

\begin{document}
\maketitle

\begin{abstract}
Post-hoc explanations of a fitted predictive model should not be used as evidence
about observed below-trend crop-yield events unless the same model has locked-period
predictive adequacy for the same target. We audit this condition in a U.S.
crop--state--year panel with
\AuditTotalRows{} rows for Barley, Canola, Oats, and Wheat from 1990--2025.
Trends and residual scales are fit on training years only; selection uses
2012--2015 validation, and 2016--2025 remains locked until final evaluation. The
validation-selected \AuditSelectedModel{} \AuditSelectedFamily{} model has final
$R^2=\AuditSelectedRtwo{}$, RMSE \AuditSelectedRMSE{} t ha$^{-1}$, and paired
$\Delta$RMSE versus a zero-residual baseline of \AuditSelectedDeltaRMSE{}
\AuditGateACI{}. The selected model fails Module A, and Full-minus-Metadata Module B
also fails with $\Delta$RMSE \AuditGateBOneDeltaRMSE{} \AuditGateBOneCI{}.
Several pre-specified sensitivities are locally favorable, including the
longer-history population, one rolling-origin fold, and severe-tail error
components. They do not alter the primary decision because they evaluate a
different population or fold, or fail the complete event-recovery rule. We
therefore retain post-hoc explanation outputs only as descriptions of a fitted
function. In synthetic stress tests, the observable A+B+E policy retains
\SyntheticPolicySensitivity{} sensitivity but only \SyntheticPolicySpecificity{}
specificity, showing that predictive modules alone do not reliably detect omitted
confounding, leakage, or distributional violations. The contribution is a
reproducible claim-eligibility protocol with paired year-block uncertainty,
weather-value and event-recovery modules, calibration cases, and a warning that
full-series detrending is retrospective only.
\end{abstract}

\begin{IEEEkeywords}
crop yield, explainable machine learning, temporal validation, negative results,
reproducibility
\end{IEEEkeywords}

\section{Introduction}
Machine-learning yield forecasts are increasingly built from weather, remote
sensing, management, and crop-model variables. Studies report useful predictive
systems across large and small panels, diverse model families, and hybrid
formulations \cite{Paudel2021,Meroni2021,LengHall2020,Klompenburg2020}.
Yet a plausible feature attribution is not causal evidence, and agreement between
explanation methods cannot compensate for a model that does not predict the future
target reliably. This distinction is especially important for below-trend yield
events, where a compelling explanation of an individual fitted prediction can be
mistaken for an explanation of the observed loss.

Weather extremes and climate variability are materially associated with crop-yield
variation at several spatial scales \cite{Lesk2016,Ray2015,SchlenkerRoberts2009,Zampieri2017,Schierhorn2021}.
Those findings motivate careful predictive evaluation; they do not license an
event-level attribution from this restricted panel. We ask a narrower question:
does a leakage-safe residual model have enough future-period predictive adequacy and stability
to permit substantive weather-feature-reliance interpretation? The answer is
evaluated before XAI is interpreted.

The paper makes four contributions. First, it formulates a study-level
claim-eligibility protocol for post-hoc explanations of predictive models. Second,
it separates overall predictive adequacy (Module A), the conditional incremental
value of weather beyond metadata (Module B), and below-trend event recovery
(Module E); Weather-only versus Metadata-only remains a representation diagnostic
(Module D). Third, it evaluates the rule in simulations with known mechanisms and
pre-specified evaluation labels, revealing both abstention behavior and
substantial false permission under structural invalidity, plus empirical
abstention and permission cases. Fourth, it
releases an audit trail that binds each module decision, interval, prediction vector, and
target definition to a reproducible artifact. Train-only preprocessing and locked
testing are protocol safeguards, not claimed as standalone novelty.

\section{Related Work}
\subsection{Yield Modeling and Detrending}
Detrending separates a local historical yield trajectory from residual variation,
but the result can depend on the detrending procedure and information available at
the evaluation date \cite{Lu2017,Meng2024}. We therefore fit each crop--state trend
and residual scale on training observations only. Ridge regression, Random Forest,
and ExtraTrees provide deliberately distinct linear and tree-based model families
\cite{HoerlKennard1970,Breiman2001,Geurts2006}; the implementations use
scikit-learn \cite{Pedregosa2011}. This is a finite pre-specified comparison, not
a search over the locked test period.

\subsection{Explanation Methods and Their Boundary}
SHAP and LIME explain a model output under their respective approximation and
reference choices \cite{LundbergLee2017,Ribeiro2016}. Permutation-style analyses,
model reliance, accumulated local effects, and conditional importance likewise
describe feature dependence or behavior of a fitted predictor
\cite{Fisher2019,ApleyZhu2020,Strobl2008}. They can be useful diagnostic artifacts.
However, none turns failed prospective predictive adequacy into evidence about an
observed yield anomaly. Prior work also shows that post-hoc explanations can be
sensitive to perturbation behavior, model randomization, and deletion/retraining
benchmarks, and that feature relevance has a causal interpretation problem unless
the target estimand and reference distribution are explicit
\cite{Adebayo2018,Hooker2019,Yeh2019,Slack2020,Janzing2020}. Those tests ask
whether explanations are faithful to model behavior; the module policy below asks
whether model behavior is eligible to support a locked-period outcome claim.

\subsection{Selective Prediction and Abstention}
Selective classification and reject-option systems trade coverage for risk:
a predictor may abstain when its risk target is not satisfied
\cite{ElYanivWiener2010,GeifmanElYaniv2017}. Our unit is different. The decision
is not a per-row class label but a study-level permission to interpret a fitted
feature family. The protocol therefore couples a predictive-adequacy module with a
paired feature-group value module and reports false permission and false abstention
on synthetic scenarios with known data-generating mechanisms and pre-specified
evaluation labels.

\subsection{Evaluation and Reproducibility}
Random splits can be inappropriate when temporal or hierarchical structure is
present \cite{Roberts2017}. Leakage can create seemingly strong scientific
machine-learning results while invalidating the intended forecast claim
\cite{KapoorNarayanan2023}. Reproducible ML also requires explicit code, data
paths, settings, and numerical evidence \cite{Pineau2021}. Accordingly, every
manuscript number is generated from an auditable run: configuration, prediction
vectors, bootstrap draws, citation check, and final PDF are all versioned outputs.

\section{Data and Protocol}
\subsection{Panel, Features, and Target}
The unit is a crop--state--year row. The processed panel contains \AuditTotalRows{}
observations from 1990--2025 for Barley, Canola, Oats, and Wheat across 12 U.S.
states: Colorado, Illinois, Iowa, Kansas, Minnesota, Montana, Nebraska, North
Dakota, Oklahoma, South Dakota, Texas, and Washington. The observed panel has 40
crop--state series, with available-year counts ranging from 10 to 36 (median 36).
Yield in t ha$^{-1}$ comes from USDA NASS Quick Stats
\cite{USDANASSQuickStats}; daily maximum/minimum temperature, precipitation, and
shortwave radiation come from NASA POWER \cite{NASAPOWER2026}. The primary design
has 35 full-season weather features. A separate sensitivity design adds early-,
middle-, and late-season proxies without redefining the target or locked population.
Because the primary features require a completed crop-season weather window, the
task is a post-season scientific audit of residual-model reliance, not a
pre-harvest forecast. Supplementary Table S2 defines each weather feature by model
name, NASA POWER field, aggregation, window, unit, spatial rule, missing handling,
and availability.

\begin{table}[t]
\centering
\caption{Data-flow and final-test eligibility reconstructed from the raw inputs.}
\label{tab:dataflow}
\scriptsize
\resizebox{\columnwidth}{!}{\begin{tabular}{lrl}
\toprule
Stage & Rows & Rule \\
\midrule
raw yield file & 1291 & source table supplied in repository \\
analysis year filter & 1257 & exclude 34 raw 1989 rows because analysis starts in 1990 \\
processed model frame & 1257 & exact reconstruction from raw yields plus NASA POWER daily input \\
final test eligible & 333 & 2016-2025 rows with at least three prior series observations \\
\bottomrule
\end{tabular}
}
\end{table}

For a test interval, each crop--state linear trend and residual scale is estimated
using only its associated training rows. A row is eligible only with at least three
training observations. The validation period begins with 144 candidate rows and
retains \AuditValidationRows{} after \AuditValidationExcluded{} pre-specified
exclusions. The locked test starts with \AuditTestCandidateRows{} candidates,
excludes \AuditTestExcluded{}, and evaluates \AuditTestRows{} rows. Future yields
cannot change a previously
constructed evaluation target; this property has a unit test and row-level audit.
The learned target is the raw train-only residual in t ha$^{-1}$:
$\hat{y}_{c,s,t}=a^{(\mathrm{train})}_{c,s}+b^{(\mathrm{train})}_{c,s}t$,
$r_{c,s,t}=y_{c,s,t}-\hat{y}_{c,s,t}$, and
$z_{c,s,t}=r_{c,s,t}/\max(\hat{\sigma}_{c,s,\mathrm{train}},\epsilon)$.
The standardized residual defines events, $\mathbb{1}[z_{c,s,t}<\tau]$, only.
This distinction prevents a pooled raw-error score from being misread as a
standardized event-recovery metric. Min-history rules of 3, 5, 8, and 10 years are
reported as sensitivity analyses because both eligibility and the fitted residual
scale can change with available history. Three observations are the minimum that
leaves one positive residual degree of freedom for estimating a training residual
scale; this is not claimed to provide a stable series-specific variance estimate.
The model matrices exclude year, yield, trend, residual, standardized residual,
event labels, predictions, history length, and residual scale; the no-shortcut
schema and overlap audit are released with the artifact.

\subsection{Selection and Baselines}
Ridge $\alpha\in\{1,10\}$ and Random Forest/ExtraTrees with 160 trees and
$\mathrm{min\_samples\_leaf}\in\{1,2\}$ are evaluated on metadata-only,
weather-only, and full representations. Stochastic fits use fixed seeds
$\{7,17,27,37,47\}$; numeric imputation/scaling, categorical encoding, and model
fitting occur within the training fold. Mean per-row seed-aggregated validation
RMSE chooses the configuration, with lexicographic tie-breaking and no locked-test
access. Seed-level RMSE standard deviations are reported so tiny differences are
not presented as deterministic wins.

The final audit compares learned predictions on identical rows to zero residual,
training-mean residual, and crop training-mean residual baselines. Prediction
vectors are hashed and pairwise maximum absolute differences are retained, so any
nearly coincident baselines are demonstrated rather than assumed. We use 2,000
deterministic calendar-year block bootstrap replicates and percentile 95\%
intervals over the ten locked calendar-year blocks. This preserves
cross-sectional dependence within a year but limits interval precision. All 2,000
bootstrap replicates were valid and retained.
We do not implement the stationary bootstrap itself; we cite it for the
dependent-resampling motivation and resample complete calendar years as paired
blocks \cite{PolitisRomano1994}. State-year and crop--state cluster schemes are
reported as pre-specified sensitivity analyses.

\begin{figure*}[t]
  \centering
  \includegraphics[height=0.40\textheight,keepaspectratio]{figure_workflow_us.pdf}
\caption{Sequential claim-eligibility workflow. Module A determines whether
claims may extend beyond model description, Module B controls weather-specific
predictive reliance, and Module E is evaluated only for an event-level claim.
Failure of Module E blocks event-recovery interpretation but retains the
weather-specific predictive-reliance claim when Modules A and B pass. Module D
is a diagnostic outside the permission path.}
  \label{fig:workflow}
\end{figure*}

\subsection{Pre-specified Claim Modules}
Let $P^\ast$ denote the locked target population implied by the split, target
construction, feature set, and evaluation period. Module A estimates
$\theta_A=\mathrm{RMSE}_{P^\ast}(f^\ast)-\mathrm{RMSE}_{P^\ast}(f_0)$, where
$f^\ast$ is the validation-selected model and $f_0$ is the zero-residual baseline.
Module B estimates
$\theta_B=\mathrm{RMSE}_{P^\ast}(f_{M+W})-\mathrm{RMSE}_{P^\ast}(f_M)$, the
incremental weather value beyond metadata. The reported estimates replace
$P^\ast$ by the locked empirical rows, use paired calendar-year block resampling,
and pass only when the upper 95\% interval is below zero.

Module E is invoked only for below-trend event claims. Its primary tail is $z<-1$;
$z<-1.5$ and $z<-2$ are sensitivities because their smaller samples reduce power.
The event module requires tail RMSE/MAE improvement, rank recovery with a positive
lower bootstrap bound and within-year permutation $p<.05$, and top-$k$ recovery
above chance by lift, hypergeometric probability, permutation $p$, and year-block
lift interval. Module D is the Weather-only versus Metadata-only representation
diagnostic and is recorded outside the permission path. The fixed policy is in
Table~\ref{tab:gate}.

\begin{table*}[t]
\centering
\caption{Fixed claim modules. Module A controls overall predictive claims,
Module B controls weather-specific predictive reliance, Module E controls
below-trend event-recovery claims, and Module D is diagnostic only. Full
configuration details are in the reproducibility manifest.}
\label{tab:gate}
\footnotesize
\renewcommand{\arraystretch}{1.12}
\begin{tabular}{p{0.16\textwidth}p{0.24\textwidth}p{0.34\textwidth}p{0.16\textwidth}}
\toprule
Module & Question & Pass condition & Claim scope \\
\midrule
A -- Overall adequacy & Does the locked model improve on the same-target baseline? & Upper paired 95\% CI of $\Delta$RMSE $<0$. & Overall predictive claim \\
B -- Weather value & Does Full improve on Metadata-only on the same locked rows? & Upper paired 95\% CI of $\Delta$RMSE $<0$. & Weather-specific predictive reliance \\
E -- Event recovery & Does the model recover the primary below-trend tail beyond null checks? & Tail RMSE/MAE, rank, permutation, and top-$k$ checks all pass. & Observed below-trend event claim \\
D -- Diagnostic & How does Weather-only compare with Metadata-only? & Reported only; not a decision module. & Descriptive diagnostic only \\
\bottomrule
\end{tabular}

\end{table*}

\begin{center}
\fbox{\begin{minipage}{0.94\linewidth}\scriptsize
\textbf{Algorithm 1: Locked claim-eligibility audit}\par
\textbf{Input:} split, candidate grid, feature groups, baseline, requested claim type.\par
1. Fit target construction and preprocessing on training rows only.\par
2. Select model/configuration on validation; lock model, feature roles, and claim rule.\par
3. Refit through the validation cutoff and predict locked future rows.\par
4. Estimate Module A on paired locked rows.\par
5. Estimate Module B with Full and Metadata-only representatives.\par
6. If an event-level claim is requested, estimate Module E on the primary tail.\par
7. Record Module D as a diagnostic outside the permission path.\par
8. Return the highest claim level whose required modules pass.
\end{minipage}}
\end{center}

\section{Results}
\subsection{Validation and Locked Test}
Table~\ref{tab:selection} shows leading validation configurations. The selected
model is \AuditSelectedModel{} with \AuditSelectedFamily{} features. The complete
grid, seed-level metrics, tie rule, and locked-test access flag are supplied in the
artifact manifest. Module A uses this single overall validation-selected
configuration. Module B is a pre-specified fixed-architecture contrast: Full and
Metadata-only predictions reuse the Module A configuration ID and swap only the
feature family. Their metrics and row/target/prediction hashes are recorded in
the selection manifest; they are not candidates for replacing Module A after the
locked period is observed.

\begin{table}[t]
\centering
\caption{Validation selection on 2012--2015. RMSE is computed after fixed-seed
prediction aggregation; Seed SD is the across-seed RMSE standard deviation, and
$P$(rank 1) is the seed-level first-rank frequency for stochastic candidates;
Ridge is deterministic and shown as N/A. Role separates the Module A winner from
the fixed-architecture Module B/D representatives.}
\label{tab:selection}
\scriptsize
\resizebox{\columnwidth}{!}{\begin{tabular}{lllrrrccc}
\toprule
Model & Features & Role & $n$ & RMSE & Seed SD & $P$(rank 1) & 1-SE & Selected \\
\midrule
ExtraTrees (leaf=1) & Weather only & Module A; D & 140 & 0.384 & 0.003 & 0.60 & Yes & Yes \\
ExtraTrees (leaf=2) & Weather only & Validation candidate & 140 & 0.385 & 0.001 & 0.20 & No & No \\
Random Forest (leaf=1) & Weather only & Validation candidate & 140 & 0.387 & 0.003 & 0.20 & No & No \\
Random Forest (leaf=2) & Weather only & Validation candidate & 140 & 0.387 & 0.002 & 0.00 & No & No \\
ExtraTrees (leaf=2) & Full & Validation candidate & 140 & 0.389 & 0.003 & 0.00 & No & No \\
Random Forest (leaf=1) & Full & Validation candidate & 140 & 0.390 & 0.002 & 0.00 & No & No \\
ExtraTrees (leaf=1) & Full & Module B Full & 140 & 0.391 & 0.003 & 0.00 & No & No \\
Ridge ($\alpha=10$) & Metadata only & Best Ridge candidate & 140 & 0.407 & N/A & N/A & No & No \\
ExtraTrees (leaf=1) & Metadata only & Module B/D Metadata & 140 & 0.407 & 0.000 & 0.00 & No & No \\
\bottomrule
\end{tabular}
}
\end{table}

The selected weather-only ExtraTrees configuration ranks first in 3/5 stochastic
validation seeds. Its advantage is small relative to several seed standard
deviations, so we retain the pre-specified selection without claiming clear
superiority. The one-standard-error flag and every seed-level
score are retained; no final-test statistic changes this selection.

On the locked test, the selected model has $R^2=\AuditSelectedRtwo{}$ and RMSE
\AuditSelectedRMSE{} t ha$^{-1}$. Its paired RMSE difference from zero crosses
zero, so Module A fails. The selected configuration deteriorates from validation
RMSE 0.384 in 2012--2015 to locked-test RMSE 0.669 in 2016--2025, consistent with
the temporal instability summarized below. Table~\ref{tab:baselines} includes all
same-task naive baselines and the selected configuration across feature families.
Module B compares Full versus Metadata-only on the same locked rows. Module D
compares Weather-only and Metadata-only as an exploratory representation
diagnostic outside the permission path. All comparisons use paired predictions on
identical locked rows. Modules A and B fail; Module D shows no clear difference
and remains diagnostic.

The Metadata-only vector nearly coincides with the zero-residual baselines on the
locked rows. The target is an OLS residual constructed separately within each
crop--state series, so its training mean is approximately zero within each fitted
series. Metadata features are time-invariant context variables (latitude,
longitude, region, and crop) and provide no within-series temporal signal. In the
locked ExtraTrees configuration, the maximum absolute prediction difference is
$1.32\times10^{-13}$ versus the zero-residual vector, $1.13\times10^{-13}$ versus
the training-mean baseline, and $1.05\times10^{-13}$ versus the crop-training-mean
baseline. Consequently, Module D is numerically identical to Module A in this panel,
although the contrasts answer different conceptual questions. Module B remains
distinct because the Full model can use weather jointly with metadata.

\begin{table*}[t]
\centering
\caption{Locked 2016--2025 same-task audit. Learned-model predictions are seed
aggregated. $\Delta$RMSE is model minus zero-residual RMSE with paired 95\%
year-block intervals; RMSE is t ha$^{-1}$.}
\label{tab:baselines}
\scriptsize
\resizebox{\textwidth}{!}{\input{../../../paper/generated/table_final_baselines.tex}}
\end{table*}

\subsection{Event-Detection Diagnostics}
Event detection is evaluated over all 333 locked rows, not only a tail preselected
by the observed outcome. The score is negative predicted residual and the score
threshold is fixed at zero before final evaluation. Event scores retain weak
global discrimination: ROC-AUC is 0.641 for the primary $z<-1$ tail and 0.672 for
the two severe-tail sensitivities; primary-tail PR-AUC is 0.338 versus prevalence
0.219. This does not contradict the module failure. ROC-AUC summarizes pairwise
ordering over all locked rows, whereas Module E asks whether the model improves
paired error and recovers the pre-specified tail beyond rank and top-$k$ null
checks. The available discrimination is diagnostically positive but insufficient
for the requested observed-event claim.

\subsection{Below-Trend Event Module and Robustness}
The broad below-trend subset contains \AuditTailRows{} locked-test rows. Table
\ref{tab:tail} separates error, rank-null, and chance-adjusted event-recovery
components. The primary $z<-1$ top-10 lift is below random expectation and its
permutation result is not significant; rank uncertainty also crosses zero. Module E
therefore fails without relying on a directional-sign statistic. Because Modules A
and B fail, neither the local nor global explanation outputs can be used to
support observed-event or weather-specific predictive claims in the primary
analysis.

\begin{table*}[t]
\centering
\caption{Tail components on train-only targets. Panel A reports error and rank
recovery; Panel B reports chance-adjusted Top-10 recovery. $H$ is the
hypergeometric probability and $P$ is the within-year permutation probability.
The Top-10 component passes only when lift $>1$, the lower lift CI $>1$, and both
$H$ and $P$ are $<.05$. Only $z<-1$ is primary; $z<-1.5$ passes error but fails
rank and top-k, while $z<-2$ fails the complete event module.}
\label{tab:tail}
\footnotesize
\renewcommand{\arraystretch}{1.08}
\textbf{Panel A: Error and rank recovery}\par\smallskip
\input{../../../paper/generated/table_tail_error_rank.tex}
\par\medskip
\textbf{Panel B: Chance-adjusted Top-10 recovery}\par\smallskip
\begin{tabular}{llrrrrrrl}
\toprule
Role & Tail & $n$ & Observed & Expected & Lift [95\% CI] & $H$ $p$ & $P$ $p$ & Status \\
\midrule
Primary & z$<$-1 & 73 & 1/10 & 1.37 & 0.73 [0.00, 2.82] & 0.794 & 0.793 & FAIL \\
Sensitivity & z$<$-1.5 & 52 & 1/10 & 1.92 & 0.52 [0.37, 2.19] & 0.907 & 0.935 & FAIL \\
Sensitivity & z$<$-2 & 36 & 1/10 & 2.78 & 0.36 [0.32, 1.68] & 0.979 & 0.985 & FAIL \\
\bottomrule
\end{tabular}

\end{table*}

\begin{table*}[t]
\centering
\caption{Locked module decisions. Module A is required for overall predictive
claims, Module B for weather-specific predictive reliance, Module E for
below-trend event-recovery claims, and Module D is diagnostic only.}
\label{tab:gateab}
\scriptsize
\resizebox{\textwidth}{!}{\begin{tabular}{p{0.04\textwidth}p{0.18\textwidth}p{0.24\textwidth}p{0.22\textwidth}p{0.12\textwidth}p{0.13\textwidth}}
\toprule
Mod. & Role & Contrast/check & Estimate or primary statistic & Status & Decision role \\
\midrule
A & Overall predictive adequacy & Selected Weather-only vs. zero residual & -0.005 [-0.019, 0.009] & FAIL & Blocks overall and higher claim levels \\
B & Incremental weather value & Full vs. Metadata-only & -0.012 [-0.029, 0.002] & FAIL & Blocks weather-specific predictive reliance \\
E & Event-recovery adequacy & $z<-1$ tail error, rank, and top-10 null checks & RMSE -0.043 [-0.074, 0.007]; rank 0.180; top-10 1/10 & FAIL & Error CI, rank, and top-10 checks do not all pass \\
D & Representation diagnostic & Weather-only vs. Metadata-only & -0.005 [-0.019, 0.009] & DESCRIPTIVE ONLY & CI crosses zero; diagnostic only \\
\bottomrule
\end{tabular}
}
\end{table*}

\begin{figure*}[t]
  \centering
  \includegraphics[width=0.98\textwidth]{figure_xai_abstention.pdf}
  \caption{Agricultural explanation output under abstention. SHAP family shares
  and the local signed group decomposition describe the locked fitted
  Weather-only model, while Modules A, B, and E fail on locked outcomes. These
  are fitted-function diagnostics, not event-recovery evidence.}
  \label{fig:xaiabstention}
\end{figure*}

The locked Weather-only model still produces coherent fitted-function
descriptions: grouped SHAP mass is concentrated in weather features, and the
local decomposition in Fig.~2 assigns one below-trend row to
recognizable weather groups. These outputs describe how the fitted predictor
varies with its inputs. They are therefore reported as model diagnostics rather
than evidence about observed yield losses.

The robustness artifact gives exact temporal pass counts and summary ranges for
the pre-specified stage-feature, crop, spatial, and target-protocol checks. These
checks do not reverse the central decision; they show its stability boundary rather
than optimize around it.

\subsection{Target, Cluster, and Model-Family Sensitivity}
The three-observation rule is the pre-specified technical minimum, not an
empirically optimal or stable variance threshold. Eligibility is audited at each validation/final
cutoff. Histories 8 and 10 retain 313 rather than 333 final rows, so their favorable
intervals are sensitivity results from a different population and do not alter the
primary Module A decision. The audited mechanism is structural: Canola--Oklahoma has seven
pre-2016 observations and Canola--Washington has five, so both series are removed
by the eight- and ten-year filters; all retained series have at least ten
pre-2016 observations. Thus the History 8 and 10 eligible row IDs and target
vectors are identical; their predictions differ only at floating-point precision
(maximum absolute difference $1.11\times10^{-16}$), so the displayed metrics
coincide. The artifacts separate target/model checks from resampling, so row
counts and cluster counts cannot be confused.
The retained sensitivity rows are \mbox{History 3}, History 5, History 8,
History 10, Standardized target, Huber detrending, HistGradientBoosting, and
ElasticNet.

Full target/model and paired-resampling tables are retained as supplementary
tables; paths and hashes are listed in the reproducibility manifest.

The standardized-target model is actually refit on the train-only residual $z$,
rather than obtained by rescaling raw-target predictions after the fact. It also
changes the point-estimate direction: the raw-residual model is slightly better
than zero at the point estimate ($\Delta$RMSE $=-0.005$), whereas the
standardized-target model is slightly worse ($\Delta$RMSE $=0.030$). Its paired
interval still crosses zero. Test-period $z$ values need not have unit variance:
each crop--state series uses its training-only residual scale, and a small training
scale can yield larger future standardized residuals without leakage. Replacing ordinary least-squares detrending by a
training-only Huber trend changes the primary event count from 73 to 40, but its
paired interval still crosses zero. These analyses are deliberately reported as
target-construction sensitivities, not candidate routes to a favorable result.

\subsection{Temporal Stability and Capacity}
The four rolling-origin folds use non-overlapping future blocks after 2003, 2006,
2009, and 2012 training cutoffs. Three of four folds have negative point
estimates, but only the 2013--2015 fold has a 95\% CI fully below zero. In
contrast, the 2007--2009 fold significantly favors the zero-residual baseline
($\Delta$RMSE $=0.017$, 95\% CI $[0.003,0.028]$). A prefix learning curve retrains the same validation-locked
configuration using progressively longer historical prefixes, while retaining only
crop--state series with at least three prefix observations. None of the four
prefix intervals establishes improvement over zero. This makes the conclusion
conditional on the available panel and model family, not a claim that more data
cannot help.

The exact rolling-origin and prefix learning-curve table is retained as a
supplementary table with paths and hashes in the reproducibility manifest.

\subsection{Audit Traceability and Interpretation Limits}
The E1--E10 manifest maps each numerical claim family to a concrete artifact:
E1--E4 cover splits, train-only targets, validation, resampling, predictions, and
tail metrics; E5--E10 cover feature timing, crop/spatial checks, target
construction, baselines, and tables. Full paths, commands, and hashes are in the
released acceptance audit and E1--E10 manifest.

Interpretation limits are summarized in Discussion; none justifies relaxing the
locked rule after evaluation.

\subsection{Retrospective Target Sensitivity}
Full-series detrending uses later yields to construct earlier residuals. It is
therefore a retrospective diagnostic, not a prospective target. Relative to the
train-only construction, the $z<-1$ final-period set has Jaccard overlap 0.605
with year-block 95\% interval [0.500, 0.675]; 21 prospective events are not
retrospective events and 13 retrospective events are not prospective events. A
coherent retrospective explanation would still not
repair the prospective fidelity failure.

The rendered retrospective comparison is retained in the supplementary artifact
set and indexed by the reproducibility manifest.

\begin{figure}[t]
  \centering
  \includegraphics[width=\columnwidth]{figure_state_delta_rmse_map.pdf}
  \caption{State-level heterogeneity on the locked 2016--2025 crop panel. Colors
  show the descriptive difference in RMSE between the validation-selected
  Weather-only model and the zero-residual baseline within each state; negative
  values favor the model. Labels report the state abbreviation and the number of
  locked test rows. Gray states are outside the study panel. These subgroup
  estimates illustrate spatial heterogeneity and are not separate module
  decisions. State geometry is obtained from the U.S. Census Bureau
  \cite{CensusCartographicBoundary}; processing settings and hashes are recorded
  in the map manifest.}
  \label{fig:statemap}
\end{figure}

\section{Calibration and External Cases}
\subsection{Synthetic Benchmark}
The claim-module policy is also evaluated where the mechanism and claim label are
known before evaluation. A repeated synthetic temporal benchmark covers 14 regimes
over 30 seeds each, including correlated features, omitted confounding, temporal
and geographic shift, and leakage. GT permit labels define six permissible
regimes and are used only for scoring; policy permit is computed from observable
Module A, B, and E outputs. Ungated explanations falsely permit 240 of 240
invalid runs (100.0\%, 95\% CI [98.4,100.0]); the observable A+B+E policy falsely permits
\SyntheticPolicyFalsePermissions{} invalid runs (\SyntheticPolicyFalsePermissionRate{},
95\% CI [65.2,76.6]) while falsely abstaining on
\SyntheticPolicyFalseAbstentions{} valid runs
(\SyntheticPolicyFalseAbstentionRate{}, 95\% CI [7.3,16.5]). Its permission rate
is \SyntheticPolicyPermissionRate{}, sensitivity \SyntheticPolicySensitivity{},
and specificity \SyntheticPolicySpecificity{}.

\begin{center}
\scriptsize
\textbf{Synthetic scenario-level calibration.} GT permit is used only as the
evaluation label. Policy permit is computed solely from pre-specified modules and
observable checks; the GT label and hidden data-generating mechanism are not
inputs to the policy. A, B, E, and Policy report pass or permission rates across
30 seeds for an event-recovery claim.\par\smallskip
\input{../../../paper/generated/table_synthetic_scenario_decisions_compact.tex}
\end{center}
The observable modules cannot reliably identify all structurally invalid regimes:
omitted confounding, geographic shift, leakage, and correlated features are
permitted in 30/30 seeds, and temporal drift in 29/30. These are false
permissions, not successful abstentions.

\subsection{Agricultural External-Resolution Check}
We also implemented a separate county-level winter-wheat check rather than
assuming that a finer panel would rescue the state-level result. The locked
pre-model population contains 574 counties with at least ten training-period
observations from 2000--2025, official county yields, static county interior
points, and annual NASA POWER weather aggregates. Validation selected a
weather-only ExtraTrees residual model from Ridge and ExtraTrees metadata,
weather, and full-feature candidates. On the untouched 2022--2025 holdout, its
RMSE was 13.47 versus 13.78 bushels per acre for the zero-residual baseline, but
the paired year-block 95\% CI for the difference was $[-0.67,0.26]$. Module B,
which compares Full with Metadata-only, passes ($\Delta$RMSE $=-0.71$, 95\% CI
$[-1.23,-0.34]$), but Module A remains inconclusive. The county analysis therefore
reaches a model-descriptive-only outcome and serves as a robustness case rather
than evidence of predictive success or cross-resolution transfer.

\subsection{Cross-Domain Sanity-Check Permission Case}
The PJM analysis is a cross-domain sanity-check permission case for the core
two-stage policy rather than a transfer claim. Its target is continuous electricity
demand, so the crop-specific extreme-event module is not invoked. Public daily PJM
electricity-demand and day-ahead forecast records from EIA Form EIA-930
\cite{EIA930} are fit with Calendar-only ExtraTrees features on
January--September 2024 and evaluated on locked October--December records; Full
adds the day-ahead demand forecast. Module A compares Full with a pre-specified
train-period mean-demand naive baseline and passes: locked RMSE is 77,521 versus
335,037 MWh, with paired bootstrap $\Delta$RMSE 95\% CI
$[-296.7, -221.8]\;\times\;10^3\,\mathrm{MWh}$. Module B compares Full with Calendar-only and
also passes: Calendar-only RMSE is 227,417 MWh and the Full-minus-Calendar paired
95\% CI is $[-186.0, -119.3]\;\times\;10^3\,\mathrm{MWh}$. Because both pass, permutation
importance is released for locked PJM rows. The XAI manifest binds row IDs,
hashes, seed handling, and output scale; agricultural XAI artifacts remain
descriptive because the required agricultural modules fail. These are predictive-reliance
interpretations, not causal or transfer claims. This intentionally easy positive
case checks that the policy is not always-abstaining; synthetic regimes define the
policy boundary.

\section{Discussion}
\subsection{Claim Scope and Multiplicity}
Several analyses are locally favorable. The eight- and ten-year history filters
improve paired error on a smaller 313-row population; the 2013--2015 rolling fold
passes, whereas the 2007--2009 fold significantly favors the baseline; severe
tails improve RMSE and MAE but fail rank and top-$k$ recovery; and the county
analysis passes Module B while Module A remains inconclusive. These results
identify design conditions under which predictive signal may be stronger, but
they either change the population, represent non-primary temporal folds, or fail
another required module. We therefore use them to motivate future study design
rather than to replace the primary locked test.

The audit separates primary decisions from descriptive measurements before the
locked period is examined. Module A, primary Module B, and the primary $z<-1$
Module E define the below-trend event permission rule; more severe tails,
alternative histories, detrending choices, resampling units, stage proxies,
crop/state diagnostics, and expanded model families are sensitivity or diagnostic
evidence. They clarify robustness and design trade-offs but are not extra chances
to declare the primary claim successful.

\subsection{Reproducibility and Traceability}
The release records input hashes, target construction, feature availability,
forbidden-column checks, locked configuration, seeds, row IDs, predictions,
bootstrap draws, and table inputs. Regenerated records are matched to the audit
record, and each paired comparison can be recomputed on identical locked rows.

\subsection{Limitations and External Validity}
State-representative weather can miss within-state exposure, weather covariates
are correlated, and the panel lacks soil, irrigation, cultivar, management, and
phenology data. The primary interval resamples only ten calendar-year blocks, and
cluster sensitivities encode different dependence assumptions without removing
that limitation. This negative result limits the present target, resolution,
features, and model family; it does not establish that weather lacks scientific
relevance for yield. The corrected synthetic benchmark also shows that predictive
adequacy and feature-group value cannot identify omitted confounding, leakage,
correlated-feature structures, or unobserved distributional violations by
themselves. GT labels are used only for evaluation, not for policy decisions.
Future work needs prospectively versioned inputs, finer exposure measurements,
design-stage validity checks, and an independent future test interval.

\section{Conclusion}
This paper evaluates whether an interpretable weather-only model can support
claims about below-trend U.S. crop-yield events after validation-locked selection.
Modules A, B, and E fail for the state panel, and a separately locked county-level
winter-wheat panel remains inconclusive under Module A. Thus this fitted model
may be described, but it does not support observed-event or weather-specific
predictive-reliance claims in this panel. More generally, locked targets,
same-row baselines, paired uncertainty, and module separation define the highest
supported claim level. The corrected synthetic benchmark's
\SyntheticPolicySensitivity{} sensitivity and \SyntheticPolicySpecificity{}
specificity show that the modules are necessary for claim discipline but not
sufficient to detect structural invalidity; leakage, omitted confounding,
measurement validity, and distribution-shift checks remain design-stage
requirements.

\bibliographystyle{IEEEtran}
\bibliography{references}
\end{document}
